\documentclass[10pt,a4paper]{article}
\usepackage[a4paper,hmargin=22mm,vmargin=18mm]{geometry}
\usepackage[T1]{fontenc}
\usepackage{lmodern,amsmath,amssymb,amsthm,booktabs,array,enumitem,hyperref,microtype}
\hypersetup{colorlinks=true,linkcolor=blue,urlcolor=blue,citecolor=blue,pdftitle={Exact Outer Multiset Dimensions of Fibonacci Cubes Through Dimension Nine: An Exact Three-Landmark Ordered--Unordered Gap},pdfauthor={Ryutaro Yonezu}}
\emergencystretch=3em
\newtheorem{theorem}{Theorem}
\newtheorem{proposition}[theorem]{Proposition}
\newtheorem{lemma}[theorem]{Lemma}
\newcommand{\odim}{\operatorname{odim}}
\newcommand{\md}{\operatorname{md}}
\newcommand{\G}{\Gamma}
\newcommand{\F}{\mathcal{F}}
\newcommand{\rep}{\mathcal{M}}
\title{Exact Outer Multiset Dimensions of Fibonacci Cubes Through Dimension Nine\\[2mm]
\large An Exact Three-Landmark Ordered--Unordered Gap}
\author{Ryutaro Yonezu\\\small Independent Researcher}
\date{13 November 2026\\\small Preprint v1.0 --- not human peer-reviewed}
\begin{document}
\maketitle
\begin{abstract}
Let $\G_n$ be the Fibonacci cube induced by binary words of length $n$ with no consecutive $1$s.
We determine the outer multiset, ordinary metric, and full multiset dimensions
for $1\leq n\leq9$ by exhaustive, exact finite computation.
In particular, $\odim(\G_n)=1,1,2,3,3,4,5,5,7$, respectively.
For the ninth Fibonacci cube we obtain the ordinary metric dimension $\dim(\G_9)=4$
and both the outer and full multiset dimensions $\odim(\G_9)=\md(\G_9)=7$.
The full-multiset versus ordinary-metric gap is at most one through $n=8$,
but is exactly three at $n=9$. This comparison requires all vertices to be
distinguished, so the outer relaxation does not account for the latter gap.
The lower bound at dimension nine uses exhaustive rejection of all
$\binom{89}{6}=581{,}106{,}988$ six-landmark sets, with an independent encoding-based computation,
while the upper bound comes from an explicit seven-landmark certificate.
We also give an explicit four-landmark outer multiset resolving set of $\G_6$ that loses
its resolving property when one additional landmark is added.
All conclusions are about the stated finite instances, not a formula for arbitrary $n$.
This is an unrefereed preprint; historical novelty remains unconfirmed.
\end{abstract}
\noindent\textbf{Keywords:} Fibonacci cube; metric dimension; outer multiset dimension; multiset dimension; resolving set; computer-assisted proof; multiset distance representation.\\
\textbf{MSC (provisional):} 05C12, 05C85, 11B39.

\section{Introduction and scope}
A graph's metric dimension is defined using \emph{ordered} distance vectors to a fixed set of landmarks.
The outer multiset dimension of Gil-Pons and coauthors~\cite{GilPons2019}
instead records only the multiset of landmark distances, and demands separation only
among vertices outside the landmark set; see also recent surveys~\cite{Farhan2026,Albejani2026}.
These are different graph invariants. The Fibonacci cube is a classical induced subgraph of a hypercube,
whose vertices are counted by Fibonacci numbers; see~\cite{Klavzar2013,Cabello2011}.

The present paper gives an exactly reproducible \emph{finite classification} through dimension nine.
Its main point is the exact gap between \emph{ordered} and \emph{unordered}
landmark-distance information for $\G_9$ when \emph{all} vertices must be distinguished.
We show that the optimal seven-landmark outer-multiset witness in fact distinguishes
every vertex, including the landmarks themselves. This paper does not claim a uniform formula, a new asymptotic bound,
or priority over undiscovered or incompletely checked literature.
In particular, the \emph{Fibonacci dimension} of~\cite{Cabello2011} concerns
isometric \emph{embedding} into a Fibonacci cube, and is not the usual metric dimension.
Independent resolving sets studied in~\cite{Monica2017} are also a different parameter.
Araujo-Pardo and coauthors~\cite{Araujo2026} give maximum-order graph
constructions for specified diameter and multiset basis size, and determine the
outer multiset dimension for circulants $C_n(1,2)$. Those results do not state the
Fibonacci-cube values below, although a full audit of indirect implications remains open.
A broader literature audit, particularly of all implications of prior metric-dimension results, remains pending.

\section{Definitions and isometry}
Let $\F_n=\{x\in\{0,1\}^n:x_ix_{i+1}=0\ (1\le i<n)\}$ and
let $\G_n$ be the induced subgraph of the $n$-cube on $\F_n$, with two vertices adjacent
when their Hamming distance equals one. With $F_0=0,F_1=1$ and $F_{m+2}=F_{m+1}+F_m$,
$|V(\G_n)|=F_{n+2}$.

For $S=\{s_1,\ldots,s_k\}\subseteq V(G)$ and $v\in V(G)$, put
$\rep(v\mid S)=\{\!\{d_G(v,s):s\in S\}\!\}$, the \emph{unordered} multiset of distances.
The set $S$ is an \emph{outer multiset resolving set} if
$\rep(u\mid S)\ne\rep(v\mid S)$ for all distinct $u,v\in V(G)\setminus S$.
The minimum possible $|S|$ is $\odim(G)$.
A \emph{full multiset resolving set} distinguishes \emph{every} pair of vertices
(including those in $S$) by their unordered distance multisets; its minimum size
is denoted $\md(G)$, or infinity if no such set exists. Every full multiset resolving
set is outer multiset resolving, so $\odim(G)\le\md(G)$.
For comparison, a usual resolving set distinguishes \emph{all} vertices by the ordered vector
$(d_G(v,s_1),\ldots,d_G(v,s_k))$; its minimum size is $\dim(G)$.

\begin{lemma}[Isometry]
For any $x,y\in\F_n$, $d_{\G_n}(x,y)=d_H(x,y)$.
\end{lemma}
\begin{proof}
Every edge changes one bit, so graph distance is at least Hamming distance.
Starting from $x$, delete all $1$s in positions where $y$ has a $0$.
Every intermediate word still avoids consecutive $1$s.
Then insert, one by one, the remaining $1$s specified by $y$.
At each step the set of $1$ positions is a subset of that of $y$, so it is legal.
This path changes each differing bit exactly once, giving the reverse inequality.
\end{proof}

\section{Exact finite results}
\begin{theorem}[Exact finite resolving parameters]\label{thm:table}
For the Fibonacci cubes $\G_n$ with $1\le n\le9$,
\[
\begin{array}{c|rrrrrrrrr}
 n&1&2&3&4&5&6&7&8&9\\\hline
 |V(\G_n)|&2&3&5&8&13&21&34&55&89\\
 \dim(\G_n)&1&1&2&2&3&3&4&4&4\\
 \odim(\G_n)&1&1&2&3&3&4&5&5&7\\
 \md(\G_n)&1&1&3&3&3&4&5&5&7
\end{array}.
\]
These are finite, computer-assisted equalities, supported by independent BFS
enumeration through $n=8$ and by the separately audited complete certificates
at $n=9$; no formula for larger $n$ is asserted.
\end{theorem}

\begin{theorem}[Exact nine-dimensional separation]\label{thm:gap}
For the ninth Fibonacci cube,
\[
\dim(\G_9)=4,\qquad
\odim(\G_9)=\md(\G_9)=7,\qquad
\md(\G_9)-\dim(\G_9)=3.
\]
The lower bounds are finite exhaustive computational certificates, not general-$n$ structural bounds.
In particular, the gap of three holds between the ordinary \emph{ordered} metric dimension
and the full \emph{unordered} multiset dimension, independently of the outer relaxation.
\end{theorem}

\subsection{Explicit upper certificates}\label{sec:witness}
For $\G_8$, a full multiset resolving set of size five is
\[
S_8=\{01001010,01010001,01010010,01010101,10010101\}.
\]
The independent BFS verifier checks all $\binom{55}{5}=3{,}478{,}761$
five-element sets and finds exactly two full multiset resolving sets.
A usual four-landmark resolving set is
\[
R_8=\{00000100,00101001,01001010,10100010\}.
\]
The same verifier exhaustively excludes all smaller cardinalities for both
parameters. Consequently $\dim(\G_8)=4$ and $\md(\G_8)=5$.
Together with Theorem~\ref{thm:gap}, this exhibits a jump in the
ordered--unordered gap from one at $n=8$ to three at $n=9$.

All words below are length-nine vertices of $\G_9$.
An outer multiset resolving set of size seven is
\begin{align*}
S_9=\{&010101000,\ 010100001,\ 010101010,\ 100101010,\\
&010001001,\ 000100100,\ 010010010\}.
\end{align*}
The exact BFS-based verifier checks that the $89-7=82$ vertices outside $S_9$
have pairwise distinct sorted distance lists. A new independent verifier further checks
that \emph{all 89} vertices of $\G_9$, \emph{including the seven landmarks},
have pairwise distinct unordered distance multisets. Therefore $S_9$ is a full
multiset resolving set, and $\md(\G_9)\le7$. The previously established
$\odim(\G_9)=7$ and $\odim(G)\le\md(G)$ imply $\md(\G_9)=7$.
An ordinary (ordered) four-landmark resolving set is
\[
R_9=\{000101001,\ 010010101,\ 100100101,\ 101010001\}.
\]
Its four-component ordered distance vectors are pairwise distinct for all $89$ vertices.
The frozen witness script additionally contains one explicit outer resolving set for
\emph{every} row of Theorem~\ref{thm:table}, as well as upper certificates
$\odim(\G_{10})\le7$ and $\odim(\G_{11})\le9$.
The latter two are \emph{upper bounds only} and are not asserted to be exact.

\subsection{Exhaustive lower certificates}\label{sec:exhaust}
For $n\le6$, a frozen Python verifier independently constructs the graph and
its BFS distance matrix, and enumerates every $k$-subset for $0\le k<\odim(\G_n)$.
A new, separately written C++17 verifier constructs BFS distances for every
$1\le n\le8$, checks all pairs against Hamming distance, and exhaustively
excludes every landmark-set size below each reported value of $\dim$ and $\md$.
For $\md(\G_8)$ it checks all five-element sets, not merely until a witness
is found, and records exactly two successes. Multiset signatures are represented
in base seven (with no carries for the tested landmark sizes), unlike the frozen
16-adic implementation.
For $n=7,8,9$ the source C++17 verifier tests all $k$-subsets for $2\le k$ below
the claimed optimum; the complete counts are as follows.
\begin{center}
\begin{tabular}{@{}rrrl@{}}\toprule
$n$ & $k$ & $\binom{F_{n+2}}{k}$ & Result \\\midrule
7 & 2 & 561 & NONE \\
7 & 3 & 5,984 & NONE \\
7 & 4 & 46,376 & NONE \\
8 & 2 & 1,485 & NONE \\
8 & 3 & 26,235 & NONE \\
8 & 4 & 341,055 & NONE \\
9 & 2 & 3,916 & NONE \\
9 & 3 & 113,564 & NONE \\
9 & 4 & 2,441,626 & NONE \\
9 & 5 & 41,507,642 & NONE \\
9 & 6 & 581,106,988 & NONE \\\bottomrule
\end{tabular}
\end{center}
The cases $k=0$ are impossible because $\G_n$ has more than one vertex.
For $k=1$ and $n\ge3$, the exterior contains $F_{n+2}-1>n$
vertices, but a single landmark yields at most $n$ distinct positive distances.
Thus $k=1$ is also impossible for $n=7,8,9$. All smaller $k$ are therefore excluded
\emph{individually}; the outer-multiset property is not assumed to be monotone under adding landmarks.

The original C++ enumerator partitions the $k$-subsets by the first two ordered indices and
recursively lists all strictly increasing suffixes. For each subset, it excludes the
landmarks, computes the multiset-distance signature of every outside vertex,
and rejects the subset on the first repeated signature.
For $n\le9$ and $k\le6$, a distance histogram is stored in base $16$
(one four-bit counter per possible distance). No counter can exceed six,
so there are no carries; every signature has an exact, collision-free representation.
A per-subset open-addressing table is used only for checking equality of those signatures.
In every NONE case the reported number of tested subsets equals the binomial coefficient.

As a differential check, a separately implemented C++20 program encodes each six-landmark
distance histogram in \emph{base seven}, as the sum of powers $7^{d(v,s)}$.
Because all multiplicities are at most six, there are no carries in base seven.
The largest relevant sum is at most $6\cdot7^9<2^{32}$.
It independently rejects all $581{,}106{,}988$ possible six-element subsets of $V(\G_9)$.
The differing encodings reduce dependence on a single implementation of the distance multiset.

Finally, an independent C++17 verifier checks all $\binom{89}{3}=113{,}564$ ordered-distance
three-landmark subsets, finding none resolving, and explicitly verifies a four-landmark
witness. Unlike outer multiset resolution, the ordinary resolving property \emph{is}
monotone under enlarging the landmark set, so rejecting all three-element sets excludes
all smaller sizes. Together these checks prove Theorem~\ref{thm:gap} within the stated
computer-assisted proof model.

\section{Failure of upward closure of outer multiset resolving sets}
\begin{proposition}\label{prop:nonmono}
There exists an outer multiset resolving set $S$ of $\G_6$ and a proper superset
$T\supset S$ that is not outer multiset resolving.
\end{proposition}
\begin{proof}
Set
\[
S=\{001010,010001,010010,010101\},\qquad
T=S\cup\{000001\}.
\]
Direct enumeration of the $21$ vertices of $\G_6$, using the isometric Hamming metric,
shows that all $17$ vertices outside $S$ have distinct distance multisets to $S$.
However, $u=001001$ and $v=010100$ both lie outside $T$ and satisfy
\[
\rep(u\mid T)=\rep(v\mid T)=\{\!\{1,2,2,3,4\}\!\}.
\]
Thus $S$ resolves the outside vertices whereas $T$ does not.
The independently written script \texttt{verify\_nonmonotonicity.py} repeats these checks
using a freshly constructed BFS graph.
\end{proof}
This is a concrete example for the stated graph, not a priority claim that
nonmonotonicity of multiset resolving sets is new in general graph theory.

\section{Reproducibility, limitations and priority status}
The new C++17 program \texttt{independent\_bfs\_n1\_to8.cpp} proves the
additional finite entries in Theorem~\ref{thm:table} by BFS and complete lower-cardinality enumeration.
Its saved log includes the exhaustive $n=8$, $k=5$ count of $3{,}478{,}761$.
The supplementary script \texttt{verify\_md\_gamma9.py} independently constructs the
Fibonacci cube as a graph, computes distances using BFS, and verifies the full
89-vertex unordered-distance certificate; it does not use any frozen verifier code.
The public reproducibility package contains eight unchanged verification sources extracted
from the frozen research archives and the v0.4 package. These include the Python BFS
witness verifier and original C++17 exhaustive search, the independent base-seven C++20
enumeration, an independent ordered-distance metric-dimension checker, and a second BFS
verifier. The original archives are retained privately, byte-for-byte. A separate portable
runner and public manifests accompany the unchanged sources.
The author reproduced the complete checks on Windows 10 / PowerShell 5.1, including both
581,106,988-case exhaustions. The saved Windows result ZIP passed independent integrity
and SHA256 checks for all six payload files; its results agree with the independent
evidence-audit report. This is independent inspection of author-run saved evidence,
not an auditor-operated Windows execution. Raw logs and personal paths are excluded
from the public distribution.

The results do not establish any lower bound matching the supplied upper certificates
for $n=10$ or $n=11$; no formula is proposed for arbitrary dimensions.
The quantity $\dim(\G_9)=4$ is presented as a computational comparator,
\emph{not} as an independently verified new historical discovery.
Existing research on ordinary resolving parameters and
Fibonacci graph families, including~\cite{Monica2017}, has not been exhaustively
checked for all possible implications, nor have all 2026 articles been completely
screened. Hence novelty remains \textbf{INCONCLUSIVE} and references \textbf{PARTIAL}.
The supplementary multiset calculation uses the same seven-point witness but a fresh
BFS implementation; the general inequality $\odim\le\md$ supplies its lower bound.
These are AI-assisted finite computations and internal audits, not human peer review.

\paragraph{Publication status.} This is \textbf{Preprint v1.0}, released under the
author's explicit authorization after the publication gates. Independent AI audit:
\textbf{PASS within the documented scope}. Author-run Windows reproduction and independent
inspection of its saved evidence: \textbf{PASS}. Human peer review has not been performed.
Novelty remains \textbf{INCONCLUSIVE}; references remain \textbf{PARTIAL}.

\vspace{1em}
\noindent\textcopyright{} 2026 Ryutaro Yonezu. All rights reserved.

\begin{thebibliography}{9}\setlength{\itemsep}{2pt}
\bibitem{GilPons2019} R. Gil-Pons, Y. Ram\'irez-Cruz, R. Trujillo-Rasua, I. G. Yero,
\emph{Distance-based vertex identification in graphs: The outer multiset dimension},
Applied Mathematics and Computation \textbf{363} (2019), 124612.
\url{https://doi.org/10.1016/j.amc.2019.124612}.
\bibitem{Klavzar2013} S. Klav\v{z}ar,
\emph{Structure of Fibonacci cubes: A survey},
Journal of Combinatorial Optimization \textbf{25} (2013), 505--522.
\url{https://doi.org/10.1007/s10878-011-9433-z}.
\bibitem{Cabello2011} S. Cabello, D. Eppstein, S. Klav\v{z}ar,
\emph{The Fibonacci dimension of a graph},
Electronic Journal of Combinatorics \textbf{18} (2011), P55.
\url{https://doi.org/10.37236/542}.
\bibitem{Monica2017} Chris Monica M. and D. Little Femilin Jana,
\emph{Independent Resolving Number of Fibonacci Cubes and Extended Fibonacci Cubes},
International Journal of Computer Applications \textbf{165}(7) (2017), 46--48.
\url{https://doi.org/10.5120/ijca2017913936}.
\bibitem{Araujo2026} G. Araujo-Pardo, C. Dalf\'o, M. Olsen, M. A. Reyes, R. Simanjuntak,
\emph{On the outer multiset dimensions of graphs},
Bolet\'in de la Sociedad Matem\'atica Mexicana \textbf{32} (2026), article 142.
\url{https://doi.org/10.1007/s40590-026-00973-3}.
\bibitem{Farhan2026} M. Farhan, S. Klav\v{z}ar, D. Kuziak, I. G. Yero,
\emph{Multiset resolvability parameters in graphs: A survey with new results and open problems},
arXiv:2607.10311 (2026). \url{https://arxiv.org/abs/2607.10311}.
\bibitem{Albejani2026} A. Albejani, Y. Lin, J. Ryan, K. A. Sugeng,
\emph{A survey on multiset dimension and its variations},
Statistics, Optimization \& Information Computing \textbf{16}(5) (2026), 4960--4988.
\url{https://doi.org/10.19139/soic-2310-5070-4062}.
\end{thebibliography}
\end{document}
